\section{Equivalence semantics and the frequency-first dictionary (protocol \texorpdfstring{$\Leftrightarrow$}{<->} physics)}
\label{app:equivalence_semantics}

This section upgrades a recurring narrative claim in the HPA--$\Omega$ program into an audit-facing contract:
physical-language statements are to be read as statements about explicit mathematical objects, \emph{modulo declared equivalence relations}, and all additional closure choices are to be performed by CAP on explicit finite candidate families.
It introduces \emph{no new axioms} beyond the two declared primitives:
tick as the executed input stream (Axiom~\ref{ax:readout_sequentiality}) and CAP as the unique deterministic closure/selection rule (Axiom~\ref{ax:cap}).

\paragraph{Why an ``equivalence section'' is needed.}
The main text already uses several invariance doctrines (tick-origin shift, local fiber relabelings, holonomy cycle-type invariance, log-mismatch invariance under unit changes), but they are distributed across sections.
Here we collect them as a single semantic layer so that later continuous dynamical closures (Section~\ref{app:cap_continuum_action_closure} and Appendix~\ref{app:thermodynamics_from_equivalence}) can be stated as \emph{mathematical closures on equivalence classes} rather than as informal matching dictionaries.

\subsection{Physical objects as equivalence classes}
\label{subsec:equivalence_physical_objects}

\paragraph{Protocol objects.}
\MathTag At fixed window length $m$, protocol microstates are $m$-bit words in $\Omega_m=\{0,1\}^m$ and stable readout labels are stable types $w\in X_m$ (Sections~\ref{sec:hpa_readout}--\ref{sec:folding_core}).
At fixed Hilbert order $n$ on the chosen screen, locality is represented by the display graph $G_n$ (Definition~\ref{def:addressing_map_graph}).

\paragraph{Semantic contract.}
\InterfaceTag A \emph{physical object} in this paper is an equivalence class of protocol objects under declared equivalence relations that capture representational freedom (e.g.\ choice of origins, basis relabelings, coarse-graining maps).
A \emph{physical observable} is an invariant functional on these equivalence classes, or (when coarse graining is included) a functional that is monotone under the declared coarse-graining preorder.

\subsection{Minimal equivalence relations used implicitly in the main text}
\label{subsec:equivalence_relations_minimal}

We record the minimal equivalence relations that are already used (often implicitly) throughout the paper.
Each item below is a \emph{semantic quotient}: it does not add new dynamical assumptions; it fixes what is meant by ``the same physics''.

\paragraph{(E1) Tick-origin shift.}
\InterfaceTag Since protocol observables depend on tick differences, $t\sim t+t_0$ is a coordinate convention (Section~\ref{subsec:tick_time_arrow}).

\paragraph{(E2) Projection-fiber equivalence (finite observability).}
\MathTag At fixed $m$, microstates $k,k'\in\{0,\dots,2^m-1\}$ are observationally equivalent if they project to the same stable type:
$k\sim_m k'$ iff $\mathrm{Fold}_m(k)=\mathrm{Fold}_m(k')$.
This is the formal core of finite observability and is the origin of degeneracy/fiber data.
%
\begin{proposition}[Observable quotient at fixed window length]\label{prop:omega_quotient_equals_Xm}
\MathTag Fix $m\ge 1$ and let $K_m:=\{0,\dots,2^m-1\}$.
Define $k\sim_m k'$ by $\mathrm{Fold}_m(k)=\mathrm{Fold}_m(k')$.
Then $\sim_m$ is an equivalence relation and the induced map
\[
\overline{\mathrm{Fold}}_m:\ K_m/\!\sim_m\ \longrightarrow\ X_m,\qquad [k]\mapsto \mathrm{Fold}_m(k),
\]
is a bijection.
Equivalently, the stable-type set $X_m$ is the observable quotient of the microstate index set at resolution $m$.
In particular, $\mathrm{Fold}_m$ factors canonically as
\[
K_m\ \twoheadrightarrow\ K_m/\!\sim_m\ \xrightarrow[\cong]{\ \overline{\mathrm{Fold}}_m\ }\ X_m.
\]
\end{proposition}
\begin{proof}
Reflexivity, symmetry, and transitivity of $\sim_m$ are immediate from the definition.
The map $\overline{\mathrm{Fold}}_m$ is well-defined because $\mathrm{Fold}_m$ is constant on equivalence classes by construction.
It is surjective because $\mathrm{Fold}_m$ is surjective onto $X_m$ (Proposition~\ref{prop:foldm_surjective}).
It is injective because if $\overline{\mathrm{Fold}}_m([k])=\overline{\mathrm{Fold}}_m([k'])$, then $\mathrm{Fold}_m(k)=\mathrm{Fold}_m(k')$ and hence $k\sim_m k'$, i.e.\ $[k]=[k']$.
\end{proof}
\begin{remark}[Fibers and residual uncertainty]\label{rem:fold_fibers_residual_uncertainty}
\InterfaceTag Under this quotient, each stable label $w\in X_m$ represents a fiber $P(w):=\mathrm{Fold}_m^{-1}(w)\subset K_m$ of observationally indistinguishable microstates.
Finite preimage degeneracy is a direct consequence of Zeckendorf truncation at finite window length and the cutoff $k<2^m$; see Section~\ref{subsec:foldm_uplift} for the explicit tail interpretation and degeneracy structure.
\end{remark}

\paragraph{(E3) Local fiber-slot relabelings (finite gauge redundancy).}
\MathTag In the padded-fiber connection model (Section~\ref{sec:protocol_connections_holonomy}), relabeling local fiber slots at a vertex is $g_x\in S_r$ and acts by conjugation on loop products (Definition~\ref{def:s4_vertex_gauge}).
Therefore, loop holonomy is physical only through conjugacy invariants (e.g.\ cycle type; Proposition~\ref{prop:cycle_type_gauge_invariant}).

\paragraph{(E4) Coarse graining (stochastic/Markov morphisms).}
\InterfaceTag Any finite observer induces a coarse-graining preorder: one readout description is less informative if it is obtained from another by a stochastic map (a Markov morphism).
This is the semantic input behind monotonicity requirements used when selecting canonical statistical quadratic forms (cf.\ C{\v{e}}ncov uniqueness in the CAP action closure; Section~\ref{app:cap_continuum_action_closure}).
%
\subsection{Resolution refinement as a projective semantics (uplift in \texorpdfstring{$m$}{m})}
\label{subsec:resolution_projective_semantics}
\paragraph{Deterministic refinement/forgetting by prefix projections.}
\MathTag For $m\ge k\ge 1$, define the prefix projection
\[
\pi_{m\to k}:X_m\to X_k,\qquad
\pi_{m\to k}(w_1\cdots w_m):=w_1\cdots w_k.
\]
This map is well-defined because $X_m$ is a forbidden-word language: a prefix of an admissible word is admissible.
It is surjective because any $u\in X_k$ extends to $u0\cdots 0\in X_m$.
Moreover, the family is functorial: for $m\ge \ell\ge k$ one has $\pi_{m\to k}=\pi_{\ell\to k}\circ \pi_{m\to \ell}$.
Appendix~\ref{app:functorial_refinement} records this structure explicitly (at $k=6$) together with deterministic refinement multiplicities used in the SM-label uplift audits.

\paragraph{Projective physical objects (resolution-consistent families).}
\InterfaceTag A finite observer may change resolution (window length) from $m$ to $m'$.
To treat ``the same object'' across such refinements without adding a new axiom, we use a projective semantics:
a resolution-consistent object is a family of stable labels $(w_m)_{m\ge m_0}$ such that
\[
\pi_{m\to k}(w_m)=w_k\qquad\text{for all }m\ge k\ge m_0.
\]
Equivalently, objects live in the inverse limit $\varprojlim_{m} X_m$ under the prefix projections.
This is the canonical finite-resolution backbone behind the informal phrase ``uplift in $m$'' used elsewhere in the paper.

\paragraph{Relation to stochastic coarse graining (E4).}
\InterfaceTag The prefix maps above are the deterministic prototype of coarse graining: they forget the last $(m-k)$ bits.
Markov morphisms in (E4) are the stochastic generalization: they allow further loss of information beyond deterministic forgetting.
Monotonicity requirements (e.g.\ for entropy-like functionals or quadratic forms) are to be read as monotonicity with respect to this preorder, with deterministic forgetting as a special case.

\paragraph{(E5) Action equivalence (boundary terms and field redefinitions).}
\InterfaceTag When we speak of ``an action'' at the continuum closure level, the physical content is the induced equations of motion.
Accordingly, actions related by adding a boundary term, or by invertible local field redefinitions, are treated as physically equivalent.
CAP selection is applied to equivalence classes (by choosing canonical representatives under a stated tie-break) rather than to raw coordinate expressions.
%
\begin{proposition}[CAP is well-defined on equivalence classes (audit form)]\label{prop:cap_on_equiv_classes}
\AuditTag Let $\sim$ be an equivalence relation on a finite candidate family $\mathcal{C}$.
Assume the objective $J:\mathcal{C}\to\RR$ is invariant under $\sim$ (i.e.\ $c\sim c'\Rightarrow J(c)=J(c')$) and that the deterministic tie-break key $\kappa:\mathcal{C}\to\mathcal{K}$ (with a fixed total order on $\mathcal{K}$) is also invariant under $\sim$.
Then CAP selection on $\mathcal{C}$ induces a unique selected equivalence class $[c_\ast]_\sim\in \mathcal{C}/\!\sim$ and a canonical representative $c_\ast$ within that class.
Equivalently, the CAP output is representation-independent in the sense of the declared equivalence semantics.
\end{proposition}
\begin{proof}
Since $\mathcal{C}$ is finite, CAP selects the lexicographically minimal pair $(J(c),\kappa(c))$.
By invariance, both $J$ and $\kappa$ are constant on each equivalence class, so the induced map
\[
\mathcal{C}/\!\sim\ \longrightarrow\ \RR\times\mathcal{K},\qquad [c]\mapsto (J(c),\kappa(c)),
\]
is well-defined.
Minimizing on $\mathcal{C}$ is therefore equivalent to minimizing on the finite quotient set $\mathcal{C}/\!\sim$, yielding a unique class $[c_\ast]_\sim$ and hence a canonical representative $c_\ast$.
\end{proof}
\begin{remark}[Audit checklist: minimal failure point]\label{rem:cap_equiv_audit_failure}
\AuditTag If the tie-break key $\kappa$ depends on a coordinate expression rather than on an invariant of the equivalence class, then the CAP output can violate the equivalence semantics even when $J$ is invariant.
Accordingly, audit form requires that both the candidate family and the tie-break specification be stated at the invariant level (Appendix~\ref{app:cap_audit_template}).
\end{remark}

\paragraph{(E6) Morita equivalence as a declared scan-level quotient (optional).}
\AuditTag In the Weyl-pair scan language, the slope parameter $\alpha$ controls the commutation phase (Equation~\eqref{eq:weyl_pair}).
If one declares an additional equivalence relation that identifies scan algebras up to Morita equivalence (an $\mathrm{SL}_2(\ZZ)$ action on $\alpha$), then distinct microscopic scan parametrizations can be treated as the same geometry at the semantic level.
We record this as an optional extension of the equivalence semantics; it is not used as a premise in the folding or continuum-closure proofs.
See Appendix~\ref{app:morita_fourier_exchange} for the precise statement and standard references.

\subsection{Frequency as a primary derived quantity (frequency-first spine)}
\label{subsec:frequency_first_spine}

\paragraph{Motivation.}
\InterfaceTag In a tick-first ontology, the most primitive quantitative notion is \emph{counting}.
The next forced notion is a \emph{rate of change per tick}.
Frequency is the canonical rate: it is dimensionless in tick units and becomes the universal bridge to energy, mass, temperature, redshift, and delay once matching dictionaries are chosen.

\begin{definition}[Frequency from phase advance (tick units)]
\label{def:frequency_from_phase}
Let $\theta(t)$ be a phase variable taking values in a circle $\RR/2\pi\ZZ$ (or in a dyadic phase register $\ZZ_{2^p}$ embedded into $\TT$).
For $t_1\neq t_2$, define the (average) angular frequency in tick units by
\[
\omega(t_1,t_2):=\frac{\Delta\theta}{\Delta t}
\quad\text{with}\quad
\Delta t:=t_2-t_1,
\]
where $\Delta\theta$ denotes the phase increment in a chosen unwrapping convention (or in the discrete register).
\end{definition}

\begin{remark}[Operator-spectrum and DFT viewpoints (equivalent dictionaries)]
\label{rem:frequency_alt_definitions}
Frequency can also be defined as a spectral parameter:
time translation by one tick is represented by a shift operator whose eigenphases define frequencies (Weyl-pair viewpoint; Appendix~\ref{app:protocol_primitives}).
Operationally, for any tick-indexed observable $q(t)$ on a finite horizon, a discrete Fourier transform yields a finite spectrum; dominant peaks define effective frequencies.
In the present paper we treat these as equivalent interface dictionaries once a specific readout/phase convention is fixed.
\end{remark}

\subsection{Concept index: physical quantities as invariants/closures}
\label{subsec:concept_index}

Table~\ref{tab:concept_index} records an audit-facing concept map.
Each physical-language quantity is assigned (i) an invariant mathematical object, and (ii) the section(s) where it is defined/closed within the tick\,+\,CAP discipline.

\begin{longtable}{p{0.16\textwidth} p{0.28\textwidth} p{0.20\textwidth} p{0.20\textwidth}}
\caption{Concept index (frequency-first): physical quantities as invariants or CAP-closed outputs.}
\label{tab:concept_index}\\
\toprule
concept & mathematical object (invariant / closure output) & where fixed/closed & operational proxy (matching layer) \\
\midrule
\endfirsthead
\toprule
concept & mathematical object (invariant / closure output) & where fixed/closed & operational proxy (matching layer) \\
\midrule
\endhead
time & tick $t\in\ZZ$ (differences only) & Axiom~\ref{ax:readout_sequentiality}; Section~\ref{sec:tick_calculus} & laboratory clock ticks after calibration \\
phase & dyadic register $\ZZ_{2^p}$ and embedding $\e^{\iu\theta}$ & Appendix~\ref{app:protocol_primitives}; Section~\ref{subsec:z128_label} & phase readout / interferometry / complex $S$-parameters \\
frequency & $\omega=\Delta\theta/\Delta t$ in tick units (Definition~\ref{def:frequency_from_phase}) & Section~\ref{app:equivalence_semantics} & spectral peaks; clock ratios; redshift \\
space (display) & addressing map $A_n$ and graph $G_n$ & Section~\ref{sec:hilbert_addressing}; Definition~\ref{def:addressing_map_graph} & locality graph used for audits \\
distance & $d_n$ (graph metric) & Definition~\ref{def:protocol_distance} & hop count / minimal transport steps \\
velocity & $v=\Delta d/\Delta t$ (tick units) & Definition~\ref{def:protocol_velocity} & propagation rate; $c$ after calibration \\
gauge connection & edge transports modulo local relabeling & Section~\ref{sec:protocol_connections_holonomy} & Wilson/plaquette statistics \\
curvature (finite) & holonomy conjugacy invariant (cycle type) & Proposition~\ref{prop:cycle_type_gauge_invariant} & loop/plaquette statistics \\
metric (continuum) & Lorentzian metric representative $g_{\mu\nu}$ (CAP-closed family) & Section~\ref{app:cap_continuum_action_closure} & redshift/lensing/clock-rate templates \\
curvature (continuum) & $R_{\mu\nu\rho\sigma}$, $G_{\mu\nu}$ from $g_{\mu\nu}$ & Section~\ref{app:variational_field_equations} & weak-field/PPN limits; classical tests \\
gauge curvature & $F_{\mu\nu}$ (field strength / curvature of connection) & Section~\ref{app:variational_field_equations} & scattering/transport; effective couplings \\
mass/energy scale & frequency/clock ratio; $r(\mu)=\log_\varphi(\mu/m_e)$ & Section~\ref{subsec:mass_resolution_coordinate}; Appendix~\ref{app:time_mass_delay} & Compton clock; scattering delay \\
lapse / redshift & $N=\kappa_0/\kappa$ from overhead $\kappa$ & Appendix~\ref{app:time_mass_delay} & redshift; Shapiro delay templates \\
action (continuum) & CAP-selected action class $[S]$ on a finite candidate family & Section~\ref{app:cap_continuum_action_closure} & effective-field fit / coarse-grained cost \\
equations of motion & Euler--Lagrange / Einstein--Yang--Mills equations from $S$ & Section~\ref{app:variational_field_equations} & dynamical response; weak-field tests \\
stress-energy & $T_{\mu\nu}$ (including information/overhead sector) & Section~\ref{app:variational_field_equations} & energy density, pressure, fluxes \\
entropy & channel-count / coarse-grained state-count functional & Appendix~\ref{app:thermodynamics_from_equivalence} & $\log$ counts; thermodynamic entropy \\
temperature & conjugate to entropy via free-energy closure & Appendix~\ref{app:thermodynamics_from_equivalence} & $k_BT$ scale; noise/thermal spectra \\
force & response functional (e.g.\ $-\nabla$ of effective free energy/action) & Appendix~\ref{app:thermodynamics_from_equivalence}; Section~\ref{app:variational_field_equations} & acceleration; pressure/gradient forces \\
overhead proxy & $\chi=\log(\kappa/\kappa_0)$ and lapse $N=\e^{-\gamma\chi}$ & Section~\ref{app:overhead_to_gravity_closure} & redshift/time delay/clock slowdown \\
effective potential (weak field) & $\Phi=-\gamma c^2(\chi-\chi_0)$ & Section~\ref{app:overhead_to_gravity_closure} & Newtonian potential proxies \\
effective density (weak field) & $\rho_{\mathrm{eff}}\propto -\Delta\chi$ & Section~\ref{app:overhead_to_gravity_closure} & lensing/dynamical mass comparisons \\
$\chi$ reconstruction protocol & Hilbert binning $\to$ window words $\to$ folding stats $\to \chi(x)$ & Section~\ref{app:chi_reconstruction_protocol} & surveys / simulations / lab arrays \\
Born probabilities & $P_k=\Tr(\rho E_k)$ (POVM) & Appendix~\ref{app:quantum_measurement_born} & empirical outcome frequencies \\
RG / running in $r$ & $\dd g/\dd r=(\log\varphi)\beta(g)$ & Appendix~\ref{app:running_couplings_resolution_flow} & running couplings / threshold matching \\
cosmology as resolution flow & $f_{\mathrm{stab}}(m)=F_{m+2}/2^m$, $d_m=2^m/F_{m+2}$ & Appendix~\ref{app:cosmology_resolution_flow} & energy-budget fits; capacity growth \\
\bottomrule
\end{longtable}

\paragraph{Remark on scope.}
\AuditTag The table focuses on the quantities required for the frequency-first dynamical closure.
Standard external unit conventions ($\hbar,c,k_B$) are treated as matching-layer calibration inputs, not as additional primitives of the tick\,+\,CAP spine.

\subsection{Curvature as loop invariants (finite and continuum)}
\label{subsec:curvature_as_loops}

\paragraph{Finite curvature from holonomy.}
\MathTag In the finite protocol language, curvature is defined operationally by loop transport:
the plaquette holonomy $p_\square$ is a loop product of edge transports, and its conjugacy invariants (cycle type) are gauge invariant under local relabelings (Proposition~\ref{prop:cycle_type_gauge_invariant}).
This is the minimal, fully finite analogue of ``curvature is holonomy''.

\paragraph{Continuum curvature as an interface limit.}
\InterfaceTag In a continuum dictionary where a (gauge or Levi--Civita) connection one-form $A$ is available,
holonomy around an infinitesimal loop is controlled by the curvature two-form $F=\dd A+A\wedge A$.
In that dictionary, the finite-loop invariant above is interpreted as a coarse-grained proxy for curvature flux through the loop.
Section~\ref{app:variational_field_equations} records the resulting continuum field equations after CAP closes a minimal action family.

\subsection{Force as response: gradients of action and free energy}
\label{subsec:force_as_response}

\paragraph{Response definition (action).}
\InterfaceTag Once a continuum representative action $S$ is fixed (as a CAP-closed output; Section~\ref{app:cap_continuum_action_closure}),
``force'' is defined as the response of $S$ (or an effective reduced action $S_{\mathrm{eff}}$) to a displacement/boundary perturbation, e.g.
\[
F_i:=-\,\frac{\partial S_{\mathrm{eff}}}{\partial x^i},
\]
in any setting where a coordinate $x^i$ is part of the chosen continuum representative.
This is an equivalence-class notion: adding a boundary term changes $S$ but not its Euler--Lagrange equations, so ``force'' is to be read as an invariant of the equations-of-motion class rather than of a raw action expression.

\paragraph{Response definition (free energy / entropic force).}
\InterfaceTag In a thermodynamic closure, an effective free energy functional $\mathcal{F}$ is defined on coarse-grained state variables, and force is likewise a response
\[
F_i=-\frac{\partial \mathcal{F}}{\partial x^i}.
\]
When $\mathcal{F}=E-TS$ is used, this yields the standard decomposition into energetic and entropic components.
Appendix~\ref{app:thermodynamics_from_equivalence} records the corresponding CAP closure and the frequency-first thermodynamic dictionary.

\subsection{Entropy as state counting under equivalence and coarse graining}
\label{subsec:entropy_state_counting}

\paragraph{Counting viewpoint.}
\InterfaceTag Under finite observability, an observed stable label $w\in X_m$ represents a whole microstate fiber $P(w)=\mathrm{Fold}_m^{-1}(w)$.
The simplest protocol-level entropy associated with this residual uncertainty is therefore a counting entropy
\[
S_{\mathrm{fib}}(w):=\log |P(w)|,
\]
optionally scaled by $k_B$ in physical units.
More generally, any coarse-graining map induces macrostates as equivalence classes; entropy is the logarithm of the macrostate multiplicity, or of an effective channel capacity in boundary settings (Appendix~\ref{app:bh_wormholes_pointer}).

\begin{lemma}[Mean degeneracy and fiber-entropy control]\label{lem:mean_degeneracy_fiber_entropy_bounds}
\MathTag Fix $m\ge 1$ and let $K_m:=\{0,\dots,2^m-1\}$.
Let $\mathrm{Fold}_m:K_m\to X_m$ be the folding map and define fibers $P(w):=\mathrm{Fold}_m^{-1}(w)$ for $w\in X_m$.
Then:
\begin{enumerate}
\item The fibers partition $K_m$, hence
\[
\sum_{w\in X_m}|P(w)|=|K_m|=2^m.
\]
In particular the mean degeneracy satisfies
\[
d_m:=\frac{|K_m|}{|X_m|}=\frac{1}{|X_m|}\sum_{w\in X_m}|P(w)|.
\]
\item The uniform mean fiber-count entropy is controlled by $d_m$:
\[
\frac{1}{|X_m|}\sum_{w\in X_m}\log|P(w)|\ \le\ \log d_m.
\]
\item Writing the stable fraction as $f_{\mathrm{stab}}(m):=|X_m|/2^m$, one has $d_m=1/f_{\mathrm{stab}}(m)$.
\end{enumerate}
\end{lemma}
\begin{proof}
Since $\mathrm{Fold}_m$ is a map onto $X_m$, the sets $\{P(w)\}_{w\in X_m}$ form a disjoint partition of $K_m$, so summing their sizes gives $|K_m|=2^m$.
The mean-degeneracy identity is then immediate by dividing by $|X_m|$.
For the entropy bound, apply Jensen's inequality to the concave function $\log$ on the positive numbers $\{|P(w)|\}_{w\in X_m}$:
\[
\frac{1}{|X_m|}\sum_{w\in X_m}\log|P(w)|\le \log\!\left(\frac{1}{|X_m|}\sum_{w\in X_m}|P(w)|\right)=\log d_m.
\]
Finally $d_m=2^m/|X_m|=1/f_{\mathrm{stab}}(m)$ by definition.
\end{proof}

\begin{proposition}[Folding as conditional entropy: fiber entropy and relative entropy]\label{prop:folding_relative_entropy_decomposition}
\MathTag Fix $m\ge 1$ and let $N$ be uniformly distributed on $K_m=\{0,\dots,2^m-1\}$.
Let $W:=\mathrm{Fold}_m(N)\in X_m$ and write
\[
\mu_m(w):=\mathbb{P}(W=w)=\frac{|P(w)|}{2^m},\qquad w\in X_m,
\]
for the induced distribution on stable labels.
Let $u_m$ denote the uniform distribution on $X_m$.
Using natural logarithms (entropy in nats), one has:
\begin{enumerate}
\item \textbf{Fiber entropy identity:}
\[
H(N\mid W)=\mathbb{E}_{\mu_m}\!\left[\log |P(W)|\right].
\]
\item \textbf{Relative-entropy decomposition:}
\[
H(N\mid W)=\log d_m + D(\mu_m\Vert u_m),
\qquad
d_m=\frac{2^m}{|X_m|},
\]
where $D(\mu\Vert\nu)=\sum_w \mu(w)\log\!\bigl(\mu(w)/\nu(w)\bigr)$ is the Kullback--Leibler divergence.
\item \textbf{Entropy-rate consequence:}
\[
\lim_{m\to\infty}\frac{1}{m}H(N\mid W)=\log(2/\varphi).
\]
\end{enumerate}
\end{proposition}
\begin{proof}
\textbf{(1)} Since $W$ is a deterministic function of $N$, conditioning on $W=w$ restricts $N$ to the fiber $P(w)$.
Because $N$ is uniform on $K_m$ and all elements in $P(w)$ are equally likely given $W=w$, one has $H(N\mid W=w)=\log|P(w)|$.
Taking expectation over $w\sim \mu_m$ yields the identity.
\par
\textbf{(2)} Because $W$ is a function of $N$, one has $H(N,W)=H(N)$.
Therefore $H(N\mid W)=H(N)-H(W)=\log|K_m|-H(\mu_m)=m\log 2-H(\mu_m)$.
On the other hand,
\[
D(\mu_m\Vert u_m)
=\sum_{w\in X_m}\mu_m(w)\log\!\left(\frac{\mu_m(w)}{1/|X_m|}\right)
=-H(\mu_m)+\log|X_m|.
\]
Since $\log d_m=\log(2^m/|X_m|)=m\log 2-\log|X_m|$, adding gives
$\log d_m+D(\mu_m\Vert u_m)=m\log 2-H(\mu_m)=H(N\mid W)$.
\par
\textbf{(3)} From Proposition~\ref{prop:rm_entropy_gap_rate} one has $\log r_m=m\log(2/\varphi)+O(1)$ for the maximal degeneracy $r_m=\max_w|P(w)|$.
From Lemma~\ref{lem:entropy_gap_hidden_exponent_cosmo} one has $\log d_m=m\log(2/\varphi)+O(1)$.
Moreover,
\[
D(\mu_m\Vert u_m)=\mathbb{E}_{\mu_m}\!\left[\log\!\left(\frac{\mu_m(W)}{u_m(W)}\right)\right]
=\mathbb{E}_{\mu_m}\!\left[\log\!\left(\frac{|P(W)|}{d_m}\right)\right]
\le \log\!\left(\frac{r_m}{d_m}\right)=O(1),
\]
so $(1/m)D(\mu_m\Vert u_m)\to 0$.
Dividing the identity in (2) by $m$ therefore yields
$
\lim_{m\to\infty} (1/m)H(N\mid W)=\lim_{m\to\infty}(1/m)\log d_m=\log(2/\varphi).
$
\end{proof}

\AuditTag A numerical certificate of the identity $H(N\mid W)=\log d_m + D(\mu_m\Vert u_m)$
for $m\in\{6,\dots,16\}$ is recorded in Table~\ref{tab:folding_entropy_decomposition} (Appendix~\ref{app:generated_tables}),
generated by \texttt{scripts/exp\_folding\_entropy\_decomposition.py}.

\paragraph{Second-law semantics.}
\InterfaceTag A key reason ``irreversibility'' can be discussed without adding a new axiom is that the map from microscopic histories to observable records is many-to-one (Section~\ref{subsec:tick_time_arrow}):
coarse graining and stability folding discard information, so entropy in the counting sense is naturally nondecreasing along protocol time when described only at the coarse level.
Appendix~\ref{app:thermodynamics_from_equivalence} makes this monotonicity precise in the CAP closure language used throughout the paper.


